\chapterbackmatter{Solutions to Weinberg Exercises}
\ExerciseEditorialNote

\begin{enumerate}
\WeinbergSolution{1}
Put \(T=t-t_1>0\), and split a path into the classical path and a
fluctuation with vanishing endpoints.  The classical solution and its action
are
\[
 \begin{aligned}
 x_{\rm cl}(\tau)
 &=\frac{x_1\sin\omega(t-\tau)+x\sin\omega(\tau-t_1)}
         {\sin\omega T},\\
 S_{\rm cl}
 &=\frac{m\omega}{2\sin\omega T}
 \left[(x^2+x_1^2)\cos\omega T-2xx_1\right].
 \end{aligned}
\]
Because the action is quadratic, the fluctuation integral is independent of
the endpoints.  Its determinant is fixed by the free-particle limit, giving
\[
 K(x,t;x_1,t_1)=
 \left(\frac{m\omega}{2\pi i\sin\omega T}\right)^{1/2}
 \exp(iS_{\rm cl}).
\]
The square root is defined by continuation from \(T=0^+\); on crossing a
zero of \(\sin\omega T\), it acquires the usual Maslov phase.  Away from
those caustics,
\[
 \boxed{\quad
 dP=|K(x,t;x_1,t_1)|^2\,dx
 =\frac{m\omega}{2\pi|\sin\omega T|}\,dx .
 \quad}
\]
An exact position eigenstate is not normalizable, so this is a formal
conditional transition density rather than a normalized probability
distribution.  A physical initial wave packet must be convolved with \(K\);
then its probability integrates to one.  At
\(\sin\omega T=0\), the kernel is instead a phase times the appropriate
delta distribution.

\WeinbergSolution{2}
At an asymptotic time define
\[
 \widetilde\phi(\mathbf p)
 =(2\pi)^{-3/2}\int d^3x\,
 e^{-i\mathbf p\cdot\mathbf x}\phi(\mathbf x),
\qquad \omega_{\mathbf p}=\sqrt{\mathbf p^2+m^2}.
\]
Equation~\eqref{eq:9.2.8} gives the normalized vacuum wave functional up
to its field-independent constant:
\[
 \Psi_0[\phi]=\mathcal N
 \exp\left[-\frac12\int d^3p\,
 \omega_{\mathbf p}\widetilde\phi(\mathbf p)
 \widetilde\phi(-\mathbf p)\right].
\]
In the field representation
\(\Pi(\mathbf x)=-i\delta/\delta\phi(\mathbf x)\).  Taking the adjoint of
Eq.~\eqref{eq:9.2.7} at \(t=0\), and acting on \(\Psi_0\), therefore gives
\[
 \boxed{\quad
 \Psi_{\mathbf p}[\phi]
 \equiv\braket{\phi}{\mathbf p}
 =\sqrt{2\omega_{\mathbf p}}\,
 \widetilde\phi(-\mathbf p)\Psi_0[\phi].
 \quad}
\]
The result has the expected first Hermite polynomial multiplying the
Gaussian.  With
\(\braket{\mathbf p}{\mathbf q}=\delta^3(\mathbf p-\mathbf q)\), its norm
follows from the Gaussian covariance
\(\langle\widetilde\phi(\mathbf p)
\widetilde\phi(\mathbf q)\rangle_0
=\delta^3(\mathbf p+\mathbf q)/(2\omega_{\mathbf p})\).

For example, an incoming particle inserts
\(\sqrt{2\omega_{\mathbf p}}\widetilde\phi(-\mathbf p)\) in the initial boundary
functional, while an outgoing particle inserts its conjugate in the final
boundary functional.  Contracting that boundary field with a field at a
vertex \(x\) supplies
\[
 \begin{array}{c|c}
 \text{external state}&\text{factor at the attached vertex}\\ \hline
 \InKet{\mathbf p}&
 (2\pi)^{-3/2}(2\omega_{\mathbf p})^{-1/2}e^{ip\cdot x}\\[2pt]
 \OutBra{\mathbf p}&
 (2\pi)^{-3/2}(2\omega_{\mathbf p})^{-1/2}e^{-ip\cdot x}
 \end{array}
\]
with \(p^0=\omega_{\mathbf p}\).  The conjugate factors describe absorption at
the opposite boundary.  These are precisely the external-line rules; all
vacuum Gaussian factors cancel between the normalized numerator and
denominator of the path integral.

\WeinbergSolution{3}
Use the independent spatial components \(v_i\) and write
\[
 P^T_{ij}=\delta_{ij}-\frac{p_ip_j}{\mathbf p^2},
\qquad
 P^L_{ij}=\frac{p_ip_j}{\mathbf p^2},
\qquad \omega_{\mathbf p}=\sqrt{\mathbf p^2+m^2}.
\]
In momentum space the free Proca Hamiltonian is a sum of oscillators whose
momentum and coordinate quadratic forms are
\[
 A=P^T+\frac{\omega_{\mathbf p}^2}{m^2}P^L,
\qquad
 B=\omega_{\mathbf p}^2P^T+m^2P^L.
\]
For a Gaussian
\(\Psi_0[v]\propto\exp[-\tfrac12\int v_i(-\mathbf p)
\mathcal E_{ij}(\mathbf p)v_j(\mathbf p)\,d^3p]\), the functional
Schrödinger equation requires
\(\mathcal E A\mathcal E=B\).  Since the projectors commute,
\[
 \boxed{\quad
 \mathcal E_{ij}(\mathbf p)
 =\omega_{\mathbf p}P^T_{ij}
 +\frac{m^2}{\omega_{\mathbf p}}P^L_{ij}
 =\omega_{\mathbf p}\left(
 \delta_{ij}-\frac{p_ip_j}{\omega_{\mathbf p}^2}\right).
 \quad}
\]
Thus
\[
 \Psi_0[v]=\mathcal N\exp\left[-\frac12\int d^3p\,
 v_i(-\mathbf p)\mathcal E_{ij}(\mathbf p)v_j(\mathbf p)\right].
\]
The product of the initial and final vacuum functionals adds the
convergent boundary damping
\[
 -\epsilon\int dt\,d^3p\,
 v_i(-\mathbf p,t)\mathcal E_{ij}(\mathbf p)v_j(\mathbf p,t)
\]
to \(iS\), with an immaterial positive rescaling of \(\epsilon\).
After the auxiliary \(v^0\) is restored, inversion of the covariant
quadratic form yields
\[
 \boxed{\quad
 \Delta^{\mu\nu}(p)=
 \frac{\eta^{\mu\nu}+p^\mu p^\nu/m^2}
      {p^2+m^2-i\epsilon}.
 \quad}
\]
Terms of order \(\epsilon\) in the longitudinal numerator depend on how
the damping is extended off shell and vanish as distributions when
\(\epsilon\to0^+\).  The important result fixed by the vacuum is the common
Feynman displacement of the physical transverse and longitudinal poles.

\WeinbergSolution{4}
Fourier transform with \(\partial_\mu\mapsto ip_\mu\), and write the
quadratic action as
\[
 I_0=-\int\frac{d^4p}{(2\pi)^4}\,
 \bar\psi^\mu(-p)\mathcal K_{\mu\nu}(p)\psi^\nu(p),
\]
where
\[
 \mathcal K_{\mu\nu}=
 (i\gamma\cdot p+m)\eta_{\mu\nu}
 +\frac{i}{3}(\gamma_\mu p_\nu+\gamma_\nu p_\mu)
 -\frac13\gamma_\mu(i\gamma\cdot p-m)\gamma_\nu .
\]
The Grassmann Gaussian integral says that the contraction is the inverse
of this finite tensor--spinor matrix.  Direct multiplication, using
\(\{\gamma^\mu,\gamma^\nu\}=2\eta^{\mu\nu}\), gives the pole part
\[
 \boxed{\quad
 \Delta^{\mu\nu}(p)=
 \frac{-i\gamma\cdot p+m}{p^2+m^2-i\epsilon}
 \left[
 \eta^{\mu\nu}-\frac13\gamma^\mu\gamma^\nu
 +\frac{2p^\mu p^\nu}{3m^2}
 +\frac{i}{3m}
   (\gamma^\mu p^\nu-\gamma^\nu p^\mu)
 \right].
 \quad}
\]
With the canonical Grassmann variables of Section~9.5, the corresponding
\(\psi\)--canonical-momentum contraction has one additional right factor
\(\beta=i\gamma^0\); the displayed expression is the conventional
\(T\{\psi^\mu(x)\bar\psi^\nu(y)\}\) propagator.

On the mass shell the tensor in brackets projects onto the physical
spin-\(\tfrac32\) sector: after multiplication by the Dirac numerator it
is transverse and gamma-traceless,
\[
 p_\mu\Delta^{\mu\nu}\big|_{p^2=-m^2}=0,
\qquad
 \gamma_\mu\Delta^{\mu\nu}\big|_{p^2=-m^2}=0.
\]
It therefore removes the spin-\(\tfrac12\) components of the vector
spinor and leaves the four massive spin-\(\tfrac32\) polarizations.  The
vacuum boundary functional supplies the same
\(-i\epsilon\) displacement as for the Dirac propagator.
\end{enumerate}
